% Generated by scripts/tikz_reviews.py using the compiled book's labels.
\pdfbookmark[0]{Chapter 1}{comparison-chapter-1}
\ReviewFigure{fig-sampling-aliasing}
{Shannon interpolation with and without aliasing}
{figures/1-shannon/sampling-aliasing.png}
{figures/tikz/shannon-sampling-aliasing.pdf}
{The sampling theorem \ref{thm-shannon-sampling} and Poisson formula \eqref{eq-poisson-formula} connect the time and frequency domains at each of four stages.}
{The paired layouts follow the supplied shannon-interp reference images. Unit sample spacing is used. Exact signal/spectrum pairs are sinc(t/a)\textasciicircum{}4 and a B3(a omega/(2 pi)), for a=4 and a=2. The rows show the original signal, its weighted Dirac comb and replicated spectrum, the sinc kernel and ideal low-pass filter, and the reconstruction. In the aliased case the bottom time curve is the exact inverse transform of the clipped replicated spectrum; it differs between samples while retaining every sample value. Red shading tracks folded spectral mass; the dashed red curves show the originals in the reconstruction row. The same vertical scales are used within each domain.}
{shannon-sampling-aliasing}
{Complete figure}
{1.4}
\ReviewFigure{fig-review-shannon-sinc}
{The cardinal sine}
{figures/1-shannon/sinc.png}
{figures/tikz/shannon-sinc.pdf}
{The reconstruction kernel in Shannon interpolation uses the normalized cardinal sine.}
{The value at zero is the continuous extension, equal to one. Every nonzero integer is a zero. The plotted normalization is sin(pi x)/(pi x), as in the theorem.}
{shannon-sinc}
{Complete figure}
{1.5}
\ReviewFigure{fig-sampling-splines}
{Cardinal B-splines of degrees zero, one, and two}
{figures/1-shannon/spline.png}
{figures/tikz/shannon-splines.pdf}
{The chapter defines phi zero as the centered unit box and obtains each subsequent spline by convolution with that box.}
{Supports and heights follow the convolution definition exactly. The quadratic spline peaks at 3/4, so its integer translates are not an interpolating basis with the samples as coefficients. Endpoint values of the box follow the closed-interval convention in the text.}
{shannon-splines}
{Complete figure}
{1.6}
\ReviewFigure{fig-aliasing}
{Two frequencies with identical samples}
{figures/1-shannon/aliasing.png}
{figures/tikz/shannon-aliasing.pdf}
{A periodic sinusoidal example explains aliasing through discrete spectral lines.}
{For unit sample spacing, cos(pi x/2) and cos(3 pi x/2) agree at every integer. Frequencies are identified modulo 2 pi; the high-frequency pair folds to the low-frequency pair. This example is periodic and is not an L2 signal on the real line.}
{shannon-aliasing}
{Complete figure}
{1.7}
\ReviewFigure{fig-review-shannon-quantizer}
{Uniform scalar quantization}
{figures/1-shannon/quantizer.pdf}
{figures/tikz/shannon-quantizer.pdf}
{The quantizer returns an integer index; dequantization multiplies that index by the step size T.}
{The axes are the normalized input u/T and integer index v. Filled left and open right endpoints implement the stated half-open cells exactly. The red vertical segment measures quantization error at a fixed input; rescaling gives the bound T/2 after dequantization. The diagonal represents the unquantized value.}
{shannon-quantizer}
{Complete figure}
{1.8}
